\section{Introduction} 
\label{sec:introduction}

Deformable object manipulation (DOM) is a crucial area of research with broad applications, from household~\citep{maitin2010cloth,miller2011parametrized,gdoom} to industrial settings~\citep{miller2012geometric,zhu2022challenges}. To aid in algorithm development and prototyping, several DOM benchmarks ~\citep{softgym,plasticine} have been developed using deformable object simulators. However,  the high dimensional state and action spaces remain a significant challenge to DOM.

%% differentiable dynamics
Differentiable physics is a promising direction for developing control policies for deformable objects. It implements physical laws as differentiable computational graphs ~\citep{Brax, Hu2020DiffTaichi}, enabling the optimization of control policies with analytical gradients and therefore improving sample efficiency. Recent studies have shown that differentiable physics-based DOM methods can benefit greatly from this approach ~\citep{plasticine, DiSec, SHAC, ILD}.  To facilitate fair comparison and further advancement of DOM techniques, it is important to develop a general-purpose simulation platform that accommodates these methods.  

% various tasks
We present \algo, a differentiable simulation framework for deformable object manipulation (DOM). In contrast with current single-task benchmarks, \algo\ encompasses a diverse range of object types, including rope, cloth, liquid, and elasto-plastic materials. The platform includes tasks with varying levels of difficulty and well-defined reward functions, such as liquid pouring, cloth folding, and rope wiping. These tasks are designed to support high-level macro actions and low-level controls, enabling comprehensive evaluation of DOM algorithms with different action spaces.

% various methods
We benchmark eight competitive DOM methods across different algorithmic paradigms, including sampling-based planning, reinforcement learning (RL), and imitation learning (IL).
For planning methods, we consider model predictive control with the Cross Entropy Method (CEM-MPC) ~\citep{Richards2005RobustCM}, differentiable model predictive control~\citep{Hu2020DiffTaichi}, and a combination of the two. For RL domains, we consider Proximal Policy Optimization (PPO)~\citep{PPO} with non-differentiable dynamics, and Short-Horizon Actor-Critic (SHAC)~\citep{SHAC} and Analytic Policy Gradients (APG)~\citep{Brax} with analytical gradients. For IL, Transporter networks~\citep{transporter} with non-differentiable dynamics and Imitation Learning via Differentiable physics (ILD)~\citep{ILD} are compared. Our experiments compare algorithms with and without analytic gradients on each task, providing insights into the benefits and challenges of differentiable-physics-based DOM methods.

% implementation and usage.
\algo\ provides a deformable object simulator, \simulator, which combines recent advances in deformable object simulation algorithms~\citep{SHAC, ILD} with the high-performance computational framework JAX~\citep{jax2018github}. This integration allows for efficient auto-differentiation and parallelization across multiple accelerators, such as multiple GPUs. All task environments are wrapped with the OpenAI Gym API~\citep{OpenAIGym}, enabling seamless integration with DOM algorithms for fast and easy development. By providing a comprehensive and standardized simulation framework, we aim to facilitate algorithm development and advance the state of the art in DOM. Moreover, our experimental results show that the dynamics model  in \algo\ enables direct sim-to-real transfer to a real robot for rope manipulation, indicating the potential applicability of our simulation platform to real-world problems.